% paper/sections/12d_interface_pipeline.tex
%
% Tier III: 界面処理 primitives の検証
%   U4 = Ridge-Eikonal reinit |∇ψ|=1 復元収束
%   U5 = Smoothed Heaviside / δ_ε モーメント精度
% ═══════════════════════════════════════════════════════════════════════════════
\subsection{Tier III：界面処理 primitives}
\label{sec:verify_interface}

Tier I/II の基盤演算子の上に，界面処理パイプライン
（CLS A→F 6-stage：移流→再初期化→物性更新→曲率→延長→PPE 入力）が構築される．
本 Tier では，このパイプラインの 2 つの key primitive を単体検証する：
U4 = Ridge--Eikonal reinit による $|\nabla\psi| = 1$ 復元能力
（§~\ref{sec:U4_ridge_eikonal_reinit}），
U5 = smoothed Heaviside と $\delta_\varepsilon$ のモーメント精度
（§~\ref{sec:U5_smoothed_heaviside_delta}）．
他の primitive（CLS 移流・物性補間・曲率）は U1 (FCCD) ・U7 (BF static droplet) で
network 的に網羅される（§~\ref{sec:numerical_integrity} トレーサビリティ表）．

% ─────────────────────────────────────────────────────────────────────────────
\subsubsection{U4：Ridge--Eikonal reinit $|\nabla\psi| = 1$ 復元収束}
\label{sec:U4_ridge_eikonal_reinit}
% ─────────────────────────────────────────────────────────────────────────────

\noindent\textbf{目的：}
§~\ref{sec:ridge_eikonal} の Ridge--Eikonal 再初期化が，
偏った初期 $\psi$ 場から $|\nabla\psi| = 1$ を満たす符号付き距離プロファイルを
反復で復元する能力を検証する．

\noindent\textbf{Part 2 対応 primitive：} \#7 (Ridge--Eikonal reinit)．

\noindent\textbf{下流連関：}
RT instability,rising bubble,Galilean A/B（界面の長時間追従に必須）．

\noindent\textbf{手順とパラメタ：}
\begin{itemize}[noitemsep]
  \item 球状界面（円形 2D：$R = 0.25$，中心 $(0.5, 0.5)$，$N = 128$）
  \item 偏った初期 $\psi$ ：$\psi_0 = H_\varepsilon(\phi_{\mathrm{SDF}}/2)$
        （勾配が真値の半分）
  \item 反復回数 $n_\tau = 1, 5, 10, 20, 50$ で
        界面近傍 $|\phi| < 3h$ における $\||\nabla\psi| - 1\|_\infty$ を計測
  \item DGR 厚み補正（§~\ref{sec:dgr_thickness}）併用 vs 非併用の比較
\end{itemize}

\noindent\textbf{結果（要約）：}
反復 5 回で $\||\nabla\psi| - 1\|_\infty$ が初期値の 1\% 以下に到達し，
20 回で機械精度近傍まで収束する．
DGR 併用により $\varepsilon_{\mathrm{eff}}/\varepsilon \approx 1.03$ が
維持され，界面厚みの過剰膨張を防ぐ．
（旧 §~\ref{sec:verify_dgr} の DGR 専用検証データ：
冪等性 $\varepsilon_{\mathrm{eff}}/\varepsilon = 1.000$（100 回適用で不変），
厚み復元 $1.4 \to 1.000$，頻度ロバスト性
（DGR なし $3.97$ vs 20 step 毎 $1.03$；99 倍改善）を
本テストに統合した．）
\label{sec:verify_dgr}%  backward compat (DGR 厚み補正)

\medskip
\noindent\textbf{再初期化界面シフトの収束（旧 §~\ref{sec:verify_reinit_shift} を統合）：}
\label{sec:verify_reinit_shift}
半径 $R = 0.25$ の円形界面に対し，
1 回再初期化後の $\psi = 0.5$ 零等値面シフト $|\Delta x_0|$ は
$\Ord{h^2}$ で収束し（$N = 32 \to 256$ で次数 $2.0$），
10 回累積比は $N = 256$ で 7.6（理想値 10）であり
逐次的な平衡接近に伴う単調減衰を示す．
質量保存モニタ（§~\ref{sec:adaptive_reinit}）と組み合わせれば
有限回の適用で系統的な界面ドリフトは生じない．

\noindent\textbf{結論：}
Ridge--Eikonal reinit は $\Ord{h^2}$ で界面を平衡へ導き，
20 回反復で $\||\nabla\psi| - 1\|_\infty$ が機械精度近傍に達する．
界面シフトは 10 回累積で $\approx 7.6$ 倍（理想値 $10$）の
\textbf{条件付き合格}（$\triangle$；DGR で補償）．

\medskip
\noindent\textbf{U4 補足：CLS 質量保存と保存型再マッピング}\\
\noindent\label{sec:verify_cls_advection}%  backward compat (CLS 移流)
\noindent\label{sec:verify_dccd_conservation}%  backward compat (DCCD 質量保存性)
DCCD 空間演算子の離散質量保存性（§~\ref{sec:dccd_mass_conservation} の理論）は，
周期 BC・壁面 BC ともに $|\sum_i \tilde{f}'_i| < 5 \times 10^{-13}$ を確認した．
CLS 移流（Zalesak / 単一渦）の質量保存誤差は，界面重み付き補正
$\psi \leftarrow \psi + (\Delta M / \sum w)\,w$, $w = 4\psi(1 - \psi)$ により
全格子で機械精度 $\Ord{10^{-15}}$ を達成．
形状誤差の支配的要因は界面厚 $\varepsilon/h$（32.6\%）＞再初期化頻度（21.4\%）であり，
DCCD 減衰寄与は $\approx 2\%$（CLS 特性波長 $\lambda_{\min} \approx 9.4h$ で $H \approx 0.98$）．
WENO5 vs DCCD の比較（旧 §~\ref{sec:verify_weno5_vs_dccd}）は paper §~\ref{sec:advection_dccd_design}
の解析議論で代替する：$L_2$ 誤差は CCD $<$ WENO5 $<$ DCCD だが
界面幅保存と離散質量保存性は DCCD のみが両立する．
\label{sec:verify_weno5_vs_dccd}%  backward compat

\medskip
\noindent\textbf{動的格子 CLS 保存型再マッピング（旧 §~\ref{sec:cls_conservation} を統合）：}\\
\noindent\label{sec:cls_conservation}%
\noindent\label{app:cls_conservation}%
$\psi$ の直接線形補間＋界面重み付き質量補正により，
格子リフレッシュ周期 $K \leq 20$ で機械精度 $\Ord{10^{-16}}$ の質量保存を達成する
（LS（SDF 経由）は全条件で約 18\%；補間が支配的誤差源）．
$K = 50$ では DCCD 移流誤差の蓄積で $3.7 \times 10^{-2}$ に劣化．
$\psi$ 直接線形補間はモノトンで $\psi \in [0, 1]$ を逸脱せず，
Gibbs 的オーバーシュートが発生しないためフラックス形式の保存的リマッピングは不要である．
NS パイプライン毎ステップ格子リビルド（旧 §~\ref{sec:verify_ns_grid_rebuild}）も
DGR 再初期化の採用で $|\Delta M|/M_0 = 6.7 \times 10^{-6}$ を達成；
この diagnose は本 Tier の単体テスト範囲を超えるため
第~\ref{sec:verification}章 の rising bubble 結合検証で網羅する．
\label{sec:verify_ns_grid_rebuild}%  backward compat

% ─────────────────────────────────────────────────────────────────────────────
\subsubsection{U5：Smoothed Heaviside / $\delta_\varepsilon$ モーメント精度}
\label{sec:U5_smoothed_heaviside_delta}
% ─────────────────────────────────────────────────────────────────────────────

\noindent\textbf{目的：}
CSF（Continuum Surface Force）モデルで用いる
smoothed Heaviside $H_\varepsilon$ および
$\delta_\varepsilon = H_\varepsilon'$ の積分モーメント精度を，
界面位置 $x_\mathrm{int}$ 既知の問題で検証する．

\noindent\textbf{Part 2 対応 primitive：} \#26 (CSF + smoothed $\delta_\varepsilon$)．

\noindent\textbf{下流連関：}
Static droplet（Laplace / 200-step），capillary wave．

\noindent\textbf{手順とパラメタ：}
\begin{itemize}[noitemsep]
  \item 1D 平面界面：$x_\mathrm{int} = 0.5$，$\varepsilon = 1.5h$，$N = 16$--$256$
  \item (a) モーメント条件：
        $\int \delta_\varepsilon\,\mathrm{d}x = 1$，
        $\int x\,\delta_\varepsilon\,\mathrm{d}x = x_\mathrm{int}$
  \item (b) $\delta_\varepsilon$ vs $h$ スケーリング：
        $\varepsilon = c h$ で $c = 0.5, 1.0, 1.5, 2.0$ の収束特性を比較
\end{itemize}

\noindent\textbf{結果：}
\begin{itemize}[noitemsep]
  \item (a) 0 次モーメント：$|\int \delta_\varepsilon - 1| < 10^{-13}$（機械精度）；
        1 次モーメント：$|\int x\delta_\varepsilon - x_\mathrm{int}|$ は
        $\Ord{h^2}$ で収束（CSF 正則化幅由来の系統誤差）．
  \item (b) $\varepsilon = 1.5h$ で $\Delta p$ 系統誤差が最小，
        $\varepsilon = 0.5h$ では数値振動，$\varepsilon = 2.0h$ では過剰平滑化が支配．
\end{itemize}

\medskip
\noindent\textbf{Young--Laplace 圧力跳び（旧 §~\ref{sec:verify_young_laplace} 由来）：}
\label{sec:verify_young_laplace}
本 sub-test の応用として，$\delta_\varepsilon$ モーメント精度が
Young--Laplace 条件 $\Delta p = \kappa/\mathit{We}$ の再現に与える影響を確認する．
円形液滴（$R = 0.25$，$\rho_l/\rho_g = 1000$，$\mathit{We} = 1$）で
解析解 $\Delta p_\mathrm{exact} = 4.0$ に対し，
全格子（$N = 32$--$128$）で相対誤差 $0.16$--$0.31\%$ を達成．
NS ソルバを介した Laplace 圧検証は U7（§~\ref{sec:U7_bf_static_droplet}）で行う．

\noindent\textbf{結論：}
$\delta_\varepsilon$ の 0 次モーメントは機械精度，1 次モーメントは $\Ord{h^2}$．
CSF 正則化が paper 全体の空間精度律速（$\Ord{h^2}$）を規定することと整合．
\textbf{合格}．
